\section{The Concept of Input and Output Kernels} \label{sec_intro}
t-distributed stochastic neighbor embedding (t-SNE) is a powerful technique for
visualizing high-dimensional data in a low-dimensional space, developed in 2008
by Laurens van der Maaten and Geoffrey Hinton (\cite{key}). Despite its wide
application, the theoretical aspects and mathematical rigor of t-SNE have
rarely been discussed thoroughly. Recently, Auffinger and Fletcher proved that
if data is sampled independently from a compact probability measure $\mu_X$,
then the output of t-SNE converges to an equilibrium measure that is also
compact (\cite{auffinger-tsne}). This is one of the very first results that
rigorously establishes the theoretical guarantees of t-SNE.

We aim to generalize the results from \cite{auffinger-tsne}. In the original
version of the t-SNE algorithm, for each pair of input or output data points, a
weight is assigned based on the distance between them. This weight is used to
construct probability distributions that will be used in the loss function. A
natural question to ask is: can we generalize the result for different weights
such that the convergence still holds and the equilibrium measure still exists?
A generalized formulation allows us to further reveal the theoretical nature of
this algorithm, greatly enhancing our understanding of t-SNE and dimension
reduction. Our work starts by generalizing the algorithm with the concept of
input and output kernels.

Let $d,s \geq 2$ with $d>s$, and $n \geq 1$. Let $X_1, \cdots, X_n \in
    \mathbb{R}^d$ be the input data points and $Y_1, \cdots Y_n \in \mathbb{R}^s$
the output data points. We define the input and output kernels as follows:
\begin{definition} [Generalized Input Kernel] \label{def:input_kernel}
    The input kernel $\ki(x,y,\sigma): \mathbb{R}^d  \times \mathbb{R}^d \to \mathbb{R}^+$ is a function of the input data points and parameter $\sigma$ such that for any $x,y \in \mathbb{R}^d$ and $\sigma > 0$,
    \begin{align}
        \ki(x,y,\sigma)=\ki(y,x,\sigma).
    \end{align}
\end{definition}
\begin{definition} [Generalized Output Kernel] \label{def:output_kernel}
    The output kernel $\ko(z,w):\mathbb{R}^s \times \mathbb{R}^s \to \mathbb{R}^+$ is a function of the output data points such that for any $x,y \in \mathbb{R}^s$,
    \begin{align}
        \ko(z,w)=\ko(w,z).
    \end{align}
\end{definition}
Clearly, by letting
\begin{align}
    \ki(x,y,\sigma)=\exp\left(-\frac{||x-y||^2}{2\sigma^2}\right), \quad \ko(z,w)=\frac{1}{1+||z-w||^2},
\end{align}
we recover the original t-SNE algorithm by using both $\ki$ and $\ko$ as the weights for the input and output data points. We can also say that the original t-SNE algorithm has a Gaussian input kernel and a t-distributed output kernel.

From now on, we will use the terminology \underline{\textbf{\textit{generalized
            t-SNE}}} when we use the generalized formulation for input and output kernels
to distinguish it from the original algorithm.

Clearly, not all choices of $\ki$ and $\ko$ will yield a valid generalized
t-SNE algorithm, such as $\ki$ that diverges as $||x-y|| \to \infty$. To prove
similar results as in \cite{auffinger-tsne}, we need to impose some conditions
to both $\ki$ and $\ko$. The conditions should provide just enough constraints
for the convergence to hold, while applying to a wide range of function
choices. Towards this end, we use the following conditions for the rest of our
work. We will provide a detailed discussion on the generality of these
conditions in section \ref{tsne:sec:generality_conditions}.

\begin{condition}[Input Kernel Conditions] \label{input_kernel_conditions}
    Let $w: \mathbb{R}_{\geq 0} \to \mathbb{R}_{\geq 0}$ be a fixed function. The input kernel has the form of
    \begin{align} \label{tsne:input_kernel}
        \ki(x,x',\sigma) = \exp(-w(\sigma\norm{x-x'}^{\theta}))
    \end{align}
    Here, $\theta>0$ is a fixed constant. $w$ satisfies:

    1. $\forall t \geq 0$, $w'(t)>0$. $\lim_{t \to \infty} w(t) = \infty$.

    2. $w'(t)+tw''(t) \geq 0$ for all $t \geq 0$.

    3. The following integral converges:
    \begin{align}
        \int_0^{\infty} t^{d-1} w(t^{\theta}) \exp(-w(t^{\theta}))dt < \infty.
    \end{align}
\end{condition}
Since the all probability masses are normalized, we can assume without loss of generality that $w(0)=0$ from now on.
\begin{condition}[Output Kernel Conditions] \label{output_kernel_conditions}
    The output kernel depends only on the distance between two output data points $y$ and $y'$. In particular, $\ko(y,y') = k (\norm{y-y'})$ and that:

    1. $k: \mathbb{R}_{\geq 0} \to \mathbb{R}^+$ is decreasing and bounded.

    2. $\ko$ satisfies
    \begin{align}
        \iint_{(0,\infty)^2} \ko(y,y') dy dy' < \infty.
    \end{align}

    3. $k$ has bounded derivative.

    4. $k'(0) = 0$.
\end{condition}

\section{Setup and Notations} \label{tsne:sec:setup}
After defining the input and output kernels, we are able to provide an
introduction to the generalized t-SNE algorithm incorporating the new
terminology. Recall that the input of the generalized t-SNE consists of vectors
$X_1, \cdots, X_n \in \mathbb{R}^d$, with the purpose of finding vectors $Y_1,
    \cdots Y_n \in \mathbb{R}^s$ that serve as a good representation of the
original vectors. The algorithm starts by defining the following probability
masses:
\begin{align}
    p_{j|i} = \frac{\ki(X_i,X_j,\sigma_i)}{\sum_{k \neq i} \ki(X_i,X_k,\sigma_i)}.
\end{align}
$\sigma_i$ satisfies that for each $i = 1, \dots, n$,
\begin{align}
    -\sum_{j = 1, j \neq i}^n p_{j|i} \log p_{j|i} = \log \textbf{Perp} =: \log (n \rho).
\end{align}
$\textbf{Perp}$ is called the perplexity and $\rho \in (0,1)$ is a given constant. The purpose of $\sigma_i$ is to make sure that for each $i$, the probability distribution induced by $p_{j|i}$ has the same entropy across all $i$. Clearly, $p_{j|i}$ is not symmetric. We then define $p_{ii} = 0$ and for $1 \leq i \neq j \leq n$,
\begin{align}
    p_{ij} = \frac{p_{j|i} + p_{i|j}}{2n} = \frac{1}{2n}\left(\frac{\ki(X_i,X_j,\sigma_i)}{\sum_{k \neq i} \ki(X_i,X_k,\sigma_i)}+\frac{\ki(X_i,X_j,\sigma_j)}{\sum_{k \neq j} \ki(X_j,X_k,\sigma_j)}\right).
\end{align}
Obviously,
\begin{align}
    \sum_{\substack{j = 1 \\ j \neq i}}^n p_{j|i} = 1, \quad \sum_{1 \leq i \neq j \leq n} p_{ij} = 1.
\end{align}
As for the output vectors $Y_1, \cdots Y_n \in \mathbb{R}^s$, set for $1 \leq i \neq j \leq n$,
\begin{align}
    q_{ij} = \frac{\ko(Yi,Yj)}{\sum_{k \neq l}\ko(Y_k,Y_l)}.
\end{align}
The goal of the generalized t-SNE is to find $Y$ that minimizes the relative entropy
\begin{align}
    L_{n,\rho}(X,Y) = \sum_{1 \leq i \neq j \leq n} p_{ij} \log \left(\frac{p_{ij}}{q_{ij}}\right).
\end{align}
More precisely, the output of t-SNE $Y^*$ is defined as
\begin{align}
    Y^* = \argmin_{Y \in (\mathbb{R}^s)^n}L_n(X,Y).
\end{align}
In the expression of $L_n(X,Y)$, all $p_{ij}$ are given and $q_{ij}$ is a function of $Y_i$ and $Y_j$. Therefore, $Y^*$ can be found simply by using gradient descent. This completes the description of the algorithm itself. Now we introduce necessary notations for convergence and asymptotic analysis. Let $\mathcal{M}(\mathbb{R}^d)$ be the space of Borel probability measures on $\mathbb{R}^d$ and $\mu_X \in \mathcal{M}(\mathbb{R}^d)$ be a probability measure. Suppose all input $X = (X_1,\cdots,X_n)$ are i.i.d. random variables with probability measure $\mu_X$. Meanwhile, let $Y^* = (Y_1,\cdots,Y_n)$ be the output of the t-SNE algorithm. Consider the following empirical measure:
\begin{align}
    \mu_n := \frac{1}{n} \sum_{i=1}^n \delta(X_i,Y_i) \in \mathcal{M}(\mathbb{R}^d \times \mathbb{R}^s).
\end{align}
We are interested in its limiting measure as $n \to \infty$. To finish the setup, there are a few more remaining quantities we would like to introduce. Let $\mu \in \mathcal{M}(\mathbb{R}^d \times \mathbb{R}^s)$. For a given kernel $\ki(x,y)$, define $F_{\rho,\mu}(x,\sigma) : \mathbb{R}^d \times \mathbb{R}^+ \to \mathbb{R}$ as
\begin{align}
    F_{\rho,\mu}(x,\sigma) = \int \frac{\ki(x,x',\sigma)}{\int \ki(x,x',\sigma)d\mu(x')} \log \left(\frac{\ki(x,x',\sigma)}{\int \ki(x,x',\sigma)d\mu(x')}\right) d\mu(x') + \log \rho.
\end{align}
Note that the last $s$ variables in the integrands stay constant. It is not hard to see that $F_{\rho,\mu}(x,\sigma)$ is the continuous version of the entropy. Next, we extend the definition of $\sigma_i$ by defining $\sigma^*_{\rho,\mu} : \mathbb{R}^d \to \mathbb{R}$ the unique solution of the equation
\begin{align}
    F_{\rho,\mu}(x,\sigma) = 0.
\end{align}
The introduction of $\sigma^*$ aligns with the procedure of finding $\sigma_i$ that controls the entropy to be the same for all $i$. Of course, the existence and uniqueness of $\sigma^*$ is not trivial and will be discussed in section \ref{tsne:sec:existence_uniqueness}. For the symmetric probability mass, let $\psi: \mathbb{R}^d \to \mathbb{R}$ be a given function and $p_{\psi}: \mathbb{R}^d \times \mathbb{R}^d \to \mathbb{R}$ be defined by
\begin{align}
    p_{\psi}(x,x') = \frac{1}{2}\left( \frac{\ki(x,x',\psi(x))}{\int \ki(x,x',\psi(x))d\mu(x')} + \frac{\ki(x,x',\psi(x'))}{\int \ki(x,x',\psi(x'))d\mu(x)}\right).
\end{align}
For the extension on the output side, let $q: \mathbb{R}^s \times \mathbb{R}^s \to \mathbb{R}$ be given by
\begin{align}
    q(y,y') = \frac{\ko(y,y')}{\iint \ko(y,y')d\mu(y) d\mu(y')}.
\end{align}
Based on the construction of $p$ and $q$, define
\begin{align}
    I_{\rho}(\mu) = \iint p_{\sigma^*_{\rho,\mu}}(x,x') \log\left(\frac{p_{\sigma^*_{\rho,\mu}}(x,x')}{q(y,y')}\right) d\mu(x,y) d\mu(x',y').
\end{align}
Recall $\mu_X \in \mathcal{M}(\mathbb{R}^d)$ is an arbitrary given probability measure. We set $\mathcal{P}_X$ to be the set of probability measures on $\mathbb{R}^d \times \mathbb{R}^s$ whose marginal coincides with $\mu_X$. Formally speaking,
\begin{align}
    \mathcal{P}_X := \{\mu \in \mathcal{M}(\mathbb{R}^d \times \mathbb{R}^s): \text{the $\mathbb{R}^d$ marginal of $\mu$ is $\mu_X$}\}.
\end{align}
Lastly, define
\begin{align}
    \tilde{\mathcal{P}}_X := \left\{ \mu \in \mathcal{P}_X: \int (p_{\sigma^*_{\rho,\mu}}(x,x')-q(y,y'))\frac{\frac{\partial}{\partial y}\ko{(y,y')}}{\ko{(y,y')}} d\mu(x',y') = 0,\mu(x,y)\text{-a.e.}\right\}.
\end{align}
\section{Main Results} \label{tsne:sec:main_results}
Using the established setup, we are able to prove the same results as in
\cite{auffinger-tsne} with a broad class of input and output kernels. On top of
that, compared to the original paper, we no longer require the sample
probability measure $\mu_X$ to have $C^1$ density. We managed to relax the
condition to $C^0$.
\begin{theorem} \label{tsne:main_theorem_1}
    Assume $\mu_X$ has a continuous density $f(x)$ and compact support. Let $X_i \in \mathbb{R}^d$ be i.i.d random variables with probability measure $\mu_X$. If condition \ref{input_kernel_conditions} and condition \ref{output_kernel_conditions} hold, then for any $\rho \in (0,1)$,
    \begin{align}
        \lim_{n \to \infty} \inf_Y L_{n,\rho}(X,Y) = \inf_{\mu \in \tilde{\mathcal{P}}_X} I_{\rho}(\mu).
    \end{align}
    Moreover, there exist a measure $\mu^*$ and a sub-sequence $\{n_k\}$ such that $\mu_{n_k} \to \mu^*$ weakly and $I_{\rho}(\mu^*) = \inf_{\mu \in \tilde{\mathcal{P}}_X} I_{\rho}(\mu)$.
\end{theorem}
\begin{theorem} \label{tsne:main_theorem_2}
    Under the same assumptions of theorem \ref{tsne:main_theorem_1}, $\mu^*$ has compact support.
\end{theorem}
\begin{remark}
    In \cite{auffinger-tsne}, a lot of the lemmas are proved under the assumption that the sample distribution $\mu_X$ is sub-Gaussian. This can also be done for our work, if we replace condition \ref{input_kernel_conditions} with the following condition \ref{input_kernel_conditions_ext}:
    \begin{condition}[Alternative Input Kernel Conditions] \label{input_kernel_conditions_ext}
        Let $w: \mathbb{R}_{\geq 0} \to \mathbb{R}_{\geq 0}$ be some fixed function. The input kernel has the form of $\ki(x,x,\sigma) = \exp(-w(\sigma\norm{x-x'}^{\theta}))$. Here, $\sigma$ is used to control entropy (thus perplexity) and $\theta>0$ is a fixed constant. Define
        \begin{align} \label{tsne:kw_def}
            \mathcal{K}_w = \left\{\tilde{C} \in \mathbb{R}_{\geq 0} \bigg\rvert \exists \epsilon>0, \text{\normalfont s.t.} \lim_{n \to \infty} n^{-1/2} \exp \left(6 w\left(\left(\tilde{C} \sqrt{\log n} \right)^{\theta} \right) \right) = O(n^{-\epsilon}) \right\}.
        \end{align}
        Then $w$ satisfies:

        1. $\forall t \geq 0$, $w'(t)>0$. $\lim_{t \to \infty} w(t) = \infty$;

        2. $w'(t)+tw''(t) \geq 0$ for all $t \geq 0$;

        3. The following integral converges:
        \begin{align}
            \int_0^{\infty} t^{d-1} w(t^{\theta}) \exp(-w(t^{\theta}))dt < \infty;
        \end{align}

        4.$\exists M > 0$ s.t. $\forall t \geq 0, w'(t) \leq w(t)$, and $\sup \mathcal{K}_w > 0$.
    \end{condition}
    We will not include proofs for lemmas assuming condition \ref{input_kernel_conditions_ext} and sub-Gaussian samples, because such generalization will not carry onto the main theorems (neither do the proofs in \cite{auffinger-tsne}). However, the proofs mentioned in the original paper and adaptable to this context are still invaluable because they demonstrate the significant efforts to provide the theoretical guarantees of t-SNE for a wider range of sample data.
\end{remark}
\section{The Generality of Conditions and Our Contributions} \label{tsne:sec:generality_conditions}
\subsection{The Generality of Input and Output Kernels}
Since the main goal of our work is the generalize the theoretical guarantees of
the t-SNE algorithm, it is important that we do not lose comprehensiveness when
assuming a concrete form of both the input and output kernels. Therefore, both
condition \ref{input_kernel_conditions} and \ref{output_kernel_conditions}
warrant an explanation.

For the formula of the input kernel in \eqref{def:input_kernel}, we start by
realizing that $t$-SNE is about the relationship between all data points. The
insights we gain from the algorithm should be translation and rotation
invariant with respect to the coordinate system we choose. Hence, it is only
reasonable to assume that the input kernel is a function of $\norm{x-x'}$ for
each pair of input data points $x,x'$.

On the other hand, the weight we assign to each pair of data points is always
positive due to nature of probability mass. We thus lose no generality by
assuming the input kernel has an exponential form because obviously, the
exponential function is a smooth bijection from $\mathbb{R}$ to $\mathbb{R}^+$.
The same above can be said for the output kernel.

Additionally, though it is coincidental that in the original version of the
algorithm $\sigma_i$ can be viewed as the variance of a Gaussian random
variable, in reality $\sigma_i$'s are only used to control the entropy of each
probability distribution induced by $x_i$. With that being said, we only need
to introduce $\sigma$ in such a way that it is able to control the entropy for
each $x$. As we can see in section \ref{tsne:sec:existence_uniqueness}, simply
multiplying $\sigma$ by $\norm{x-x'}^{\theta}$ is enough to achieve this goal.
Moreover, we prove the convergence of $\sigma^*$ in section
\ref{tsne:sec:uniform_convergence}, providing the anchor of the following
results just as in \cite{auffinger-tsne}.

Combining all the above, we arrive at the form of the input kernel in
\eqref{tsne:input_kernel}, by adding a general function $w$ to the product
$\sigma \norm{x-x'}^{\theta}$. This gives us a kernel we can easily work
without losing any generality.

Besides the form of the input kernel, we also want to make sure that the
constraints we impose on the kernels in condition \ref{input_kernel_conditions}
and \ref{output_kernel_conditions} are also general enough for the significance
of our results. Since we frequently deal with integrals containing the input
and output kernels, their integrability is the bare minimum we should assume.
Along with integrability comes the requirement of their decay to 0 at infinity.
Apart from that, we do not add any excessive constraints to the kernels and
readers will see later in proofs that those assumptions are strictly necessary
for the results to hold.

Specifically, notice the second assumption of condition
\ref{input_kernel_conditions} is basically requiring the input kernel to decay
at least at a minimal speed by controlling its second order derivative. If we
let $w(t) = \log t$, we have $w'(t)+tw''(t) = 0$. Roughly speaking, this
implies that the main theorems hold as long as the input kernel decays at
polynomial rate or faster, which is a big step up than proving the convergence
for Gaussian kernels alone.
\subsection{Our Contributions}
Given all the prior context, we are able to summarize our
contributions as follows:

1. \textbf{A Foundation for Generalized Kernels:} We established a concrete, yet general, formulation for both input and output kernels, providing a solid foundation for subsequent theoretical analyses.

2. \textbf{Guaranteed Convergence:} We identified specific conditions on these kernels that ensure the existence of an equilibrium measure, $\mu^*$, and guarantee the convergence of the perplexity-controlled t-SNE algorithm. These conditions provide crucial insights into the essential properties of kernels required for optimal algorithmic performance.

3. \textbf{Relaxed Constraints and Novel Proofs:} Through the development of novel proof techniques, we demonstrated the existence and uniqueness of the perplexity-controlling parameter, $\sigma^*$, while simultaneously relaxing the constraints on the underlying probability measure, $\mu_X$. This significantly broadens the applicability of t-SNE and enables a more comprehensive convergence analysis.

4. \textbf{Extended Convergence Analysis:} We rigorously proved the existence of the equilibrium measure, $\mu^*$, and demonstrated that, under the specified conditions, the output of the generalized t-SNE algorithm converges to this equilibrium. This result dramatically extends the convergence analysis in \cite{auffinger-tsne} to a significantly wider class of kernels.
\section{Existence and Uniqueness of \texorpdfstring{$\sigma^*$}{sigma-star}}\label{tsne:sec:existence_uniqueness}
In this section we establish the existence and uniqueness of $\sigma^*_{\rho,\mu}(x)$. Recall that $\sigma^*_{\rho,\mu}(x)$ is the unique solution of the equation
\begin{align}
    F_{\rho,\mu}(x,\sigma) = \int \frac{\ki(x,x',\sigma)}{\int \ki(x,x',\sigma)d\mu(x')} \log \left(\frac{\ki(x,x',\sigma)}{\int \ki(x,x',\sigma)d\mu(x')}\right) d\mu(x') + \log \rho = 0
\end{align}
Intuitively speaking, $\sigma^*$ is the limit function of $\sigma_i$'s that are used to control the entropy. The purpose of this section is to prove that such limit exists and is unique for each given $x$. The existence and uniqueness of $\sigma^*$ will become the foundation of everything we will prove afterwards. By condition \ref{input_kernel_conditions}, we can write
\begin{align} \label{big_F_formula}
      & F_{\rho,\mu}(x,\sigma) \nonumber                                                                                                                                                                                                                    \\
    = & \int \frac{\exp(-w(\sigma\norm{x-x'}^{\theta}))}{\int \exp(-w(\sigma\norm{x-x'}^{\theta})) d\mu(x')} \log\left(\frac{\exp(-w(\sigma\norm{x-x'}^{\theta}))}{\int \exp(-w(\sigma\norm{x-x'}^{\theta})) d\mu(x')}\right)d\mu(x') + \log \rho \nonumber \\
    = & -\int \frac{w(\sigma\norm{x-x'}^{\theta})\exp(-w(\sigma\norm{x-x'}^{\theta}))}{\int \exp(-w(\sigma\norm{x-x'}^{\theta})) d\mu(x')} d\mu(x') \nonumber                                                                                               \\
      & \qquad \qquad \qquad \qquad - \log \left(\int \exp(-w(\sigma\norm{x-x'}^{\theta}))d\mu(x')\right)  + \log \rho.
\end{align}
We start with the lemma that gives $F$ an increasing lower bound.
\begin{lemma} \label{existence_uniqueness:existence_lemma}
    If $\mu$ has a continuous, bounded density $f(x)$, then there exists a constant $C>0$ such that
    \begin{align}  \label{big_F_lowerbound_final}
        F_{\rho,\mu}(x,\sigma) \geq C \log \sigma - \log f(x) + \log \rho.
    \end{align}
\end{lemma}
\begin{proof}
    Notice that for any integrable function $h$ on $\mathbb{R}_{\geq 0}$, we have
    \begin{align}
         & \int_{\mathbb{R}^d} h(\norm{x}) dx = \int_{\mathbb{S}^{d-1}} \int_0^{\infty} r^{d-1}h(r) dr d\nu=S_{d-1} \int_0^{\infty} r^{d-1}h(r) dr,
    \end{align}
    where $\mathbb{S}^{d-1}$ is the $(d-1)$-dimensional sphere, $d\nu$ the corresponding surface measure, and $S_{d-1}$ the volume of the sphere. By Condition \ref{input_kernel_conditions}, we know that
    \begin{align}
        \int_0^{\infty} t^{d-1} w(t^{\theta}) \exp(-w(t^{\theta}))dt < \infty.
    \end{align}
    Since $w$ is increasing and $\lim_{t \to \infty} w(t) = \infty$, it is easy to see that
    \begin{align}
        \int_0^{\infty} t^{d-1} \exp(-w(t^{\theta}))dt < \infty.
    \end{align}
    Denote
    \begin{align}
        Z_d :=\int_{\mathbb{R}^d}\exp(-w(\norm{\tau}^{\theta})) d \tau = S_{d-1} \int_0^{\infty} t^{d-1}\exp(-w(t^{\theta})) dt.
    \end{align}
    Consider
    \begin{align}
          & \int \exp(-w(\sigma\norm{x-x'}^{\theta}))d\mu(x') = \int \exp(-w(\sigma\norm{x-x'}^{\theta})) f(x')dx' \nonumber \\
        = & Z_d \sigma^{-d/\theta}\int \frac{\exp(-w(\sigma\norm{x-x'}^{\theta}))}{Z_d \sigma^{-d/\theta}} f(x')dx'.
    \end{align}
    Because $\lim_{x \to \infty} w(x) = \infty$, the integrand of the above integral converges weakly to $\delta_x f(x')$. To prove this statement, notice that
    \begin{align}
               & \int \frac{\exp(-w(\sigma\norm{x-x'}^{\theta}))}{Z_d \sigma^{-d/\theta}} f(x')dx' \nonumber                                          \\
        =      & \frac{1}{Z_d} \int \exp\left(-w \left(\norm{\sigma^{1/\theta}(x-x')}^{\theta}\right)\right) f(x') d\sigma^{1/\theta}(x-x') \nonumber \\
        \equiv & \frac{1}{Z_d} \int \exp(-w(\norm{\tau}^{\theta})) f(x-\tau \sigma^{-1/\theta}) d\tau.
    \end{align}
    Since $f$ is bounded, $\exp(-w(\norm{\tau}^{\theta})) \sup f$ is integrable. By Dominant Convergence Theorem and the continuity of $f$, we have
    \begin{align}
          & \lim_{\sigma \to \infty} \frac{1}{Z_d} \int \exp(-w(\norm{\tau}^{\theta})) f(x-\tau \sigma^{-1/\theta}) d\tau \nonumber \\
        = & \frac{1}{Z_d} \int \lim_{\sigma \to \infty} \exp(-w(\norm{\tau}^{\theta})) f(x-\tau \sigma^{-1/\theta}) d\tau \nonumber \\
        = & f(x) \frac{1}{Z_d} \int \exp(-w(\norm{\tau}^{\theta})) d\tau \nonumber                                                  \\
        = & f(x).
    \end{align}
    There exists some $\epsilon(\sigma) \to 0$ as $\sigma\to \infty$ and a constant $C>0$, that
    \begin{align} \label{big_f_lowerbound_1}
        -\log \left(\int \exp(-w(\sigma\norm{x-x'}^{\theta}))d\mu(x')\right) = & -\log\left(Z_d \sigma^{-d/\theta}\right) - \log f(x)+ \epsilon(\sigma) \nonumber \\
        \geq                                                                   & C\log \sigma - \log f(x).
    \end{align}
    Next, since $f$ is bounded, there exists $M>0$ such that $f \leq M$.
    \begin{align} \label{big_F_lowerbound_2}
             & \int \frac{w(\sigma\norm{x-x'}^{\theta})\exp(-w(\sigma\norm{x-x'}^{\theta}))}{\int \exp(-w(\sigma\norm{x-x'}^{\theta})) d\mu(x')} d\mu(x') \nonumber                                                                  \\
        \leq & M \frac{\int w(\sigma\norm{x-x'}^{\theta})\exp(-w(\sigma\norm{x-x'}^{\theta}))dx'}{\int \exp(-w(\sigma\norm{x-x'}^{\theta}))f(x') dx'} \nonumber                                                                      \\
        =    & M \frac{\sigma^{-d/\theta}\int w(\norm{\sigma^{1/\theta}(x-x')}^{\theta})\exp(-w(\norm{\sigma^{1/\theta}(x-x')}^{\theta}))d\sigma^{1/\theta}(x-x')}{ \int \exp(-w(\sigma \norm{(x-x')}^{\theta}))f(x') dx'} \nonumber \\
        =    & M \frac{\sigma^{-d/\theta}\int w(\norm{\tau}^{\theta})\exp(-w(\norm{\tau}^{\theta}))d\tau}{ \int \exp(-w(\sigma \norm{(x-x')}^{\theta}))f(x') dx'} \nonumber                                                          \\
        % =    & M \frac{\int w(\norm{t}^{\theta})\exp(-w(\norm{t}^{\theta}))dt}{Z_d \int \exp(-w(\sigma \norm{(x-x')}^{\theta}))f(x') dx'} \nonumber                                                                                  \\
        \to  & \frac{M}{Z_d f(x)} \int w(\norm{\tau}^{\theta})\exp(-w(\norm{\tau}^{\theta}))d\tau < \infty.
    \end{align}
    The numerator converges because of condition \ref{input_kernel_conditions}. As the bound in (\ref{big_F_lowerbound_2}) does not depend on $\sigma$, we can combine (\ref{big_F_formula}), (\ref{big_f_lowerbound_1}) and (\ref{big_F_lowerbound_2}) and deduce that
    \begin{align}
        F_{\rho,\mu}(x,\sigma) \geq C \log \sigma - \log f(x) + \log \rho.
    \end{align}
\end{proof}
This lemma proves the existence of $\sigma^*_{\rho,\mu}(x)$. Indeed, by \eqref{big_F_formula}, apparently
\begin{align} \label{existence_uniqueness:limit_1}
    \lim_{\sigma \to 0} F_{\rho,\mu}(x,\sigma) = \log \rho < 0
\end{align}
since $\rho \in (0,1)$. Now with lemma \ref{existence_uniqueness:existence_lemma}, we get
\begin{align} \label{existence_uniqueness:limit_2}
    \lim_{\sigma \to \infty} F_{\rho,\mu}(x,\sigma) = \infty.
\end{align}
By continuity of $f$ and intermediate value theorem, $\sigma^*_{\rho,\mu}(x)$ exists. To prove the uniqueness, we are left to show $F_{\rho,\mu}$ is strictly increasing. The next lemma will do just that.
\begin{lemma} \label{existence_uniqueness:uniqueness_lemma}
    Given the same condition as in lemma \ref{existence_uniqueness:existence_lemma}, $F_{\rho,\mu}(x,\sigma)$ is strictly increasing in $\sigma$.
\end{lemma}
To tackle this proof, we need one additional lemma inspired by theorem 7 in \cite{liu2000inequ}.
\begin{lemma} \label{miracle_lemma_ext}
    Let $f(x),g(x),h(x)$ be continuous functions $\mathbb{R}^d \to \mathbb{R}^+$ such that for all $x,y \in \mathbb{R}^d$
    \begin{align} \label{miracle_condition_ext}
        \left(g(x)-g(y)\right)\left(\frac{f(y)}{h(y)}-\frac{f(x)}{h(x)}\right) \geq 0.
    \end{align}
    Let $\mu$ be an arbitrary Lebesgue measure and assume all integrals in this lemma converge. Then the inequality
    \begin{align} \label{miracle_inequ_ext}
        \frac{\int_{\mathbb{R}^d} f(x)d\mu}{\int_{\mathbb{R}^d} h(x)d\mu} \geq \frac{\int_{\mathbb{R}^d} f(x)g(x)d\mu}{\int_{\mathbb{R}^d} h(x)g(x)d\mu}
    \end{align}
    holds. If (\ref{miracle_condition_ext}) reverses, then (\ref{miracle_inequ_ext}) reverses.
\end{lemma}
\begin{proof}
    The proof of this lemma is similar to the one in \cite{liu2000inequ}. It suffices to show that
    \begin{align}
        \int_{\mathbb{R}^d} f(x)d\mu \int_{\mathbb{R}^d} h(x)g(x)d\mu \geq \int_{\mathbb{R}^d} h(x)d\mu \int_{\mathbb{R}^d} f(x)g(x)d\mu.
    \end{align}
    The above inequality is equivalent to
    \begin{align}
        \int_{\mathbb{R}^d} f(x)d\mu(x) \int_{\mathbb{R}^d} h(y)g(y)d\mu(y) \geq \int_{\mathbb{R}^d} h(x)d\mu(x) \int_{\mathbb{R}^d} f(y)g(y)d\mu(y).
    \end{align}
    That means we need to show
    \begin{align}
        \mathcal{G}:=\int_{\mathbb{R}^d} \int_{\mathbb{R}^d} g(y)(f(x)h(y)-f(y)h(x))d\mu(x)d\mu(y) \geq 0.
    \end{align}
    Swap $x$ and $y$ in the formula of $\mathcal{G}$ and we have an equivalent expression:
    \begin{align}
        \mathcal{G}=\int_{\mathbb{R}^d} \int_{\mathbb{R}^d} g(x)(f(y)h(x)-f(x)h(y))d\mu(x)d\mu(y).
    \end{align}
    Adding the last two formulas together yields
    \begin{align}
        2\mathcal{G} = & \int_{\mathbb{R}^d} \int_{\mathbb{R}^d} (g(x)-g(y))(f(y)h(x)-f(x)h(y))d\mu(x)d\mu(y) \nonumber                                              \\
        =              & \int_{\mathbb{R}^d} \int_{\mathbb{R}^d} h(x)h(y)\left(g(x)-g(y)\right)\left(\frac{f(y)}{h(y)}-\frac{f(x)}{h(x)}\right)d\mu(x)d\mu(y) \geq 0
    \end{align}
    by the hypothesis of the lemma. This finishes the proof.
\end{proof}
\begin{proof}[Proof of lemma \ref{existence_uniqueness:uniqueness_lemma}]
    We compute the derivative of $F_{\rho,\mu}(x,\sigma)$ directly using \eqref{big_F_formula}. The following simplified
    notations will be used in the first several steps:
    \begin{align}
        w = w(\sigma\norm{x-x'}^{\theta}), \quad w'=w'(\sigma\norm{x-x'}^{\theta}), \quad \exp(\cdot) = \exp(-w(\sigma\norm{x-x'}^{\theta})).
    \end{align}
    Omit $d\mu(x')$ in the integration for cleaner writing, and we have
    \begin{align}
          & \frac{\partial}{\partial \sigma} F_{\rho,\mu}(x,\sigma) \nonumber                                                                                                                                    \\
        = & -\int \frac{w'\norm{x-x'}^{\theta}\exp(\cdot)}{\int \exp(\cdot)} -\int \frac{w\exp(\cdot)'}{\int \exp(\cdot)} +
        \int \frac{w \exp(\cdot) \left(\int \exp(\cdot)\right)'}{\left(\int \exp(\cdot)\right)^2} - \frac{\left(\int \exp(\cdot)\right)'}{\int \exp(\cdot)} \nonumber                                            \\
        = & -\int \frac{w\exp(\cdot)'}{\int \exp(\cdot)} +
        \int \frac{w \exp(\cdot) \left(\int \exp(\cdot)\right)'}{\left(\int \exp(\cdot)\right)^2} \quad \text{\normalfont (The first and fourth terms cancel.)} \nonumber                                        \\
        = & \int \frac{w\cdot w' \cdot \norm{x-x'}^{\theta} \cdot \exp(\cdot)}{\int \exp(\cdot)}-\frac{\int w \cdot \exp(\cdot) \int w'\cdot \norm{x-x'}^{\theta} \exp(\cdot)}{\left(\int \exp(\cdot)\right)^2}.
    \end{align}
    We are able to move differentiation inside the integral because of dominant convergence theorem. To show that $\frac{\partial}{\partial \sigma} F_{\rho,\mu}(x,\sigma) \geq 0$, it suffices to show
    \begin{align} \label{monotone_inequ}
        \int w\cdot w' \cdot \norm{x-x'}^{\theta} \cdot \exp(\cdot) \int \exp(\cdot) \geq \int w \cdot \exp(\cdot) \int w'\cdot \norm{x-x'}^{\theta} \exp(\cdot).
    \end{align}
    For each $x,y,x' \in \mathbb{R}$ if the equality
    \begin{align} \label{monotone_condition}
         & (w'(\sigma\norm{x-x'}^{\theta})\norm{x-x'}^{\theta}-w'(\sigma\norm{y-x'}^{\theta})\norm{y-x'}^{\theta}) \nonumber                                   \\
         & \qquad \qquad \qquad \qquad \qquad \qquad \cdot \left(\frac{1}{w(\sigma\norm{y-x'}^{\theta})}-\frac{1}{w(\sigma\norm{x-x'}^{\theta})}\right) \geq 0
    \end{align}
    is true, then setting $f(x) = \exp(\cdot)$, $h(x) = w \cdot \exp(\cdot)$ and $g(x) = w'\cdot \norm{x-x'}^{\theta}$ in lemma \ref{miracle_lemma_ext} shows that (\ref{monotone_inequ}) holds. Let $\norm{x-x'}^{\theta} > \norm{y-x'}^{\theta}$ without loss of generality. Then by monotonicity of $w$, we have $w(\sigma\norm{x-x'}^{\theta})>w(\sigma\norm{y-x'}^{\theta})$. In order for (\ref{monotone_condition}) to be true, we are left to show
    \begin{align}
        w'(\sigma\norm{x-x'}^{\theta})\norm{x-x'}^{\theta}-w'(\sigma\norm{y-x'}^{\theta})\norm{y-x'}^{\theta} \geq 0.
    \end{align}
    Since $\sigma >0$, the above inequality is equivalent to:
    \begin{align} \label{generalized_condition_inequ}
        w'(\sigma\norm{x-x'}^{\theta})\sigma \norm{x-x'}^{\theta}-w'(\sigma\norm{y-x'}^{\theta}) \sigma \norm{y-x'}^{\theta} \geq 0.
    \end{align}
    Let $\tilde{w}(x) = x w'(x)$. Based on the second assumption in Condition \ref{input_kernel_conditions} and mean value theorem, $\tilde{w}$ satisfies for all $a,b\geq 0$, $a>b$, there exists $c \in [b,a]$ that
    \begin{align}
        \frac{aw'(a) - bw'(b)}{a-b} = w'(c) + cw''(c) \geq 0.
    \end{align}
    This proves \eqref{generalized_condition_inequ} by letting $a = \sigma\norm{x-x'}^{\theta}$ and $b = \sigma\norm{y-x'}^{\theta}$, as desired.
\end{proof}
We can now formally state and prove the main result of this section.
\begin{proposition}
    Let $\mu \in \mathcal{M}(\mathbb{R}^d)$ with $C^0$ density $f$ such that $f$ is bounded. The equation $F_{\rho,\mu}(x,\sigma) = 0$ has a unique solution $\sigma^*_{\rho,\mu}(x)$ for all $\rho \in (0,1)$ and $x \in \mathbb{R}^d$. Moreover, $x \mapsto \sigma^*_{\rho,\mu}(x)$ is smooth.
\end{proposition}
\begin{proof}
    The existence and uniqueness can be proven by combining \eqref{existence_uniqueness:limit_1}, \eqref{existence_uniqueness:limit_2}, lemma \ref{existence_uniqueness:existence_lemma} and lemma \ref{existence_uniqueness:uniqueness_lemma}. The smoothness of $\sigma^*_{\rho,\mu}(x)$ follows from the implicit function theorem.
\end{proof}

\section{Uniform Convergence of \texorpdfstring{$F_{\rho, \mu_n}(x,\sigma)$}{F}}\label{tsne:sec:uniform_convergence}